% ============================================================
% Chapter 4 — Experimental Methodology (target 8–9 pp)
% Status: FULL DRAFT 2026-07-14.
% Sources: CLAUDE.md <measurement_rules>/<gate_policy>/<experimental_state>
%   (repo-tracked project record), docs/EXPERIMENT_REGISTRY.md (App. A),
%   episode_classification.py L69-85 (stuck/timeout split),
%   carla_multi_agent_env.py L522-565 (path_efficiency).
% All development-phase numbers below are quoted from the project registry
% and are tied to run IDs; per-candidate detail condensed in Appendix A.
% ============================================================
\chapter{Experimental Methodology}
\label{ch:methodology}

The platform of \cref{ch:system} was not designed in one pass: it is the
product of an iterative campaign in which every change had to pass a
predefined quantitative gate before entering the trunk. This chapter
documents the measurement rules (\cref{sec:metrics}), the gate protocol
(\cref{sec:gate}), the validity defects found and fixed along the way
(\cref{sec:bugs}), a summary of the iterations themselves
(\cref{sec:iterations}), and the design of the final comparison campaign
(\cref{sec:campaign}).

% ------------------------------------------------------------
\section{Metrics and Measurement Rules}
\label{sec:metrics}

\paragraph{Unit of measurement.}
The unit of analysis is the \emph{agent-level record}: one JSON entry per
agent per episode (six per episode, \cref{sec:termination}). All
aggregates are computed over records after deduplication by episode and
agent identifier (in evaluation, the scenario is part of the key, since
episode identifiers can repeat across scenarios). Two integrity conditions
are enforced before any number is reported: exactly six records per
episode, and no non-finite value in any field. A run failing either
condition is not analysed further.

\paragraph{Strict success criterion.}
An agent is successful if and only if its termination reason is
\texttt{route\_complete}. In particular, \texttt{route\_short} --- reaching
the end of a degenerately short route --- is \emph{not} success
(\cref{sec:termination}). This rule was stress-tested during development:
a change that counted \texttt{route\_short} as success was introduced and
subsequently reverted (\cref{sec:iterations}), keeping the criterion
strict at the cost of lower absolute numbers. \Cref{ch:discussion}
discusses why mid-range absolute success rates are a feature, not a
defect, of the resulting measurement: away from floor and ceiling, the
metric retains room to move in both directions, which is precisely what a
comparative study needs.

\paragraph{Three reporting views.}
Every result is reported for vehicles and pedestrians separately, plus a
combined view; the two classes have different tasks, different action
spaces and different structural limits (\cref{sec:ped-limits}), and the
experimental record shows regime effects that exist for one class only.
Joint episode-level success (``all six agents succeed'') is never used:
with correlated agents in a shared world it compounds failures
multiplicatively and hides class-specific structure.

\paragraph{Metric catalogue.}
\Cref{tab:metrics} defines the metrics used throughout
\cref{ch:results}. Cumulative values over all training episodes are the
primary aggregation; sliding-window success (window 50) is used only by
the curriculum gate (\cref{sec:regimes}); within-run dynamics are shown
by splitting training into quartiles (Q1--Q4).

\begin{table}[t]
  \centering\small
  \caption{Metric catalogue. All rates are shares of agent-level records.}
  \label{tab:metrics}
  \begin{tabular}{@{}p{3.6cm}p{9.6cm}@{}}
    \toprule
    Metric & Definition \\
    \midrule
    Success rate (SR) & records with \texttt{route\_complete} \\
    Collision rate & records with \texttt{collision} \\
    Offroad rate & records with \texttt{offroad} (vehicles) \\
    Stuck rate & truncated records with route completion $< 0.3$ or an active loop detector \\
    Timeout rate & truncated records with route completion $\geq 0.3$ \\
    Stuck$+$timeout & sum of the previous two; the primary anti-stall target \\
    Route completion & fraction of route waypoints consumed, $[0,1]$ \\
    Path efficiency & completed optimal length / actual distance travelled, capped at 1; defined as 0 for stuck records \\
    Mean speed & km/h for vehicles, m/s for pedestrians \\
    No-progress steps & steps since the last waypoint advance, at episode end \\
    Route diagnostics & \texttt{route\_source} shares, under-target and too-short flags, crossing counts (\cref{sec:routes}) \\
    \bottomrule
  \end{tabular}
\end{table}

% ------------------------------------------------------------
\section{Gate-Driven Development Protocol}
\label{sec:gate}

Run-to-run variance in this environment is large (quantified below at
$\sim$10 percentage points of vehicle SR between identical replicates),
and the available compute budget did not allow replicating every candidate
change. The protocol adopted in response has four rules.

\paragraph{1. One knob at a time.}
Every candidate is a narrow, behaviourally isolated diff --- a single
reward term, a single threshold, a single observation group. Candidates
touching the same file are evaluated sequentially, never stacked. When a
bundled change failed its gate, attribution between its halves was treated
as unresolved rather than guessed.

\paragraph{2. Paired A/B at a fixed diagnostic budget.}
Each candidate runs for $3\times10^{5}$ steps against the current trunk
baseline, with the same experiment seed and, thanks to the deterministic
route seeding of \cref{sec:seeding}, a route sequence paired with the
baseline by construction. Development-phase runs were locked to the easy
level unless the candidate targeted the curriculum itself, so that level
mixing could not confound the comparison. When runs of different lengths
are compared, records are truncated to equal episode counts.

\paragraph{3. A predefined quantitative gate.}
For vehicle-focused changes the candidate is promoted only if, versus the
paired baseline: vehicle SR improves by at least $+2.0$\,pp;
stuck$+$timeout improves by at least $-2.0$\,pp; collision and offroad
each worsen by no more than $+1.0$\,pp; and both integrity conditions of
\cref{sec:metrics} hold. Changes that touch shared machinery carry a
\emph{must-not-collapse} side condition on the other class (pedestrian SR
within $-2$\,pp). A failed gate means the candidate is reverted; the trunk
is never left in a partially evaluated state.

\paragraph{4. Mechanism and outcome are recorded separately.}
For each candidate the registry records whether the intended
\emph{mechanism} occurred (e.g.\ the critic's explained variance rising
from $\approx 0$ to $0.87$ after a value-loss fix) and whether the
\emph{behavioural gate} passed --- the two frequently dissent. A small
number of gate-failing knobs were knowingly retained with the hypothesis
recorded as falsified (e.g.\ a tenfold collision penalty that left the
collision rate flat), because reverting them would have restored a defect
or because the failure itself was informative. These exceptions are
explicit in the registry (\cref{app:registry}); nothing was retained
silently.

Within-run quartile trajectories complement the gate: a candidate whose
cumulative numbers pass but whose Q4 declines (a signature observed for
one pedestrian-pace retune) is treated as rejected.

% ------------------------------------------------------------
\section{Environment Validity: Bugs Found and Fixed}
\label{sec:bugs}

In a custom multi-agent platform the measurement apparatus is part of the
experiment. Three episodes during development produced apparent policy
effects that were in fact artifacts; they are reported here both because
they shaped the final protocol and because their detection signatures may
be useful to others.

\paragraph{Route-length bug (task distribution).}
The route planner could emit routes far longer than the level target, so
nominal ``easy'' episodes were frequently not easy. Enforcing the
$[0.5, 2.0]\times$ target bounds (\cref{sec:routes}) moved cumulative
vehicle training SR by $\approx +50$\,pp \emph{with no change to the
policy or the reward} --- a pure task-distribution change. Two
consequences were adopted: absolute metrics are only comparable within the
same environment version, and all earlier results were re-labelled as
era-specific (the vehicle floor of the April pilot,
\cref{sec:pilot}, was later traced to this defect). This episode is the
strongest internal evidence for a rule used throughout this thesis:
\emph{never compare absolute numbers across environment versions; anchor
claims to paired contrasts within a version}.

\paragraph{Route-seeding bug (broken pairing).}
Per-agent route seeds were originally derived from Python's built-in
\texttt{hash()}, which is salted per process: two nominally identical runs
drew \emph{different} route sequences, silently degrading every A/B
comparison. The fix derives seeds from an explicit
\texttt{SeedSequence} over traffic seed, reset counter and a stable digest
of the agent identifier (\cref{sec:seeding}). The lesson: pairing must be
constructed and verified (byte-identical route-generation statistics
across arms), not assumed.

\paragraph{Evaluation-pipeline bugs (measurement collapse).}
The first checkpoint evaluation returned a success rate of exactly
33.33\% in every scenario --- one of three vehicles completing, every
episode. This ``too clean to be true'' pattern exposed four coupled
defects: (i) the per-episode scenario seed was constant, so every episode
replayed the same traffic configuration; (ii) the policy was evaluated at
its deterministic mean, collapsing the Gaussian policy onto a single
trajectory per scenario; (iii) the evaluation subprocess never initialized
the episode logger, so no agent-level records were written at all; and
(iv) the torch seed was constant across episode subprocesses, so even
stochastic actions would have replayed correlated noise streams. All four
were fixed (per-episode scenario and sampler seeds, stochastic
evaluation, logger propagation with clean restart), the partial
deterministic evaluation was discarded, and the stochastic protocol of
\cref{sec:eval-pipeline} became the standard. None of these fixes touches
training; per the project's own rules, no policy claim is derived from
them.

The general lesson of this section: measurement validity has to be
established \emph{before} tuning, and implausibly regular results are a
detection signal, not good news.

% ------------------------------------------------------------
\section{Iterative Development Summary}
\label{sec:iterations}

\Cref{tab:candidates} condenses the development record; the full
per-candidate registry, with run IDs and mechanisms, is in
\cref{app:registry}. Development proceeded in six phases:

\begin{enumerate}
  \item \textbf{Architectural pilot (April 2026).} Nine runs crossing
        critic encoders and training regimes, in the pre-bugfix
        environment era. Data were integral but vehicles sat at floor SR
        (0.2--3.8\%) in every cell, making the architectural comparison
        uninformative --- and redirecting effort toward environment
        validity (\cref{sec:pilot}).
  \item \textbf{Correctness.} The route-length and route-seeding fixes of
        \cref{sec:bugs}, plus per-step diagnostic instrumentation
        (progress, no-progress steps, stuck causes, route provenance).
  \item \textbf{Observability.} Route-aware vehicle features, then the
        47D extension exposing no-progress state, loop-penalty state and
        time remaining (\cref{sec:obs-actions}), closing the
        state-aliasing gap between reward and observation. The resulting
        easy-locked baseline was healthy (SR plateau $\approx 79$\%,
        critic explained variance 0.983).
  \item \textbf{Optimizer and reward iterations.} A sequence of
        single-knob candidates with mixed verdicts (\cref{tab:candidates});
        notably, a value-function repair whose mechanism was confirmed but
        whose gate failed, and a tenfold collision penalty that left
        collisions flat --- collision avoidance was not learnable from the
        44D observations of that era, a perception limit rather than a
        reward-weight problem.
  \item \textbf{Cross-class coupling.} Faster pedestrians saturate the
        vehicles' hazard channel and can push vehicles into a defensive
        zero-speed equilibrium through the \texttt{safe\_to\_push} gate.
        The promoted resolution (V1) raised the gate threshold from 0.75
        to 0.85, keeping vehicle urgency active under moderate hazard;
        a pedestrian-pace retune in the opposite direction (V3) was
        rejected on a declining final-quartile trajectory.
  \item \textbf{Curriculum mechanics and final trunk.} Windowed unlock
        with rehearsal floors, the weakest-policy gate, and raised
        pressure caps (\cref{sec:regimes}); finally, the pedestrian
        crosswalk cap (\texttt{max\_cross} $8 \to 3$), the last change to
        enter the trunk before the campaign (vehicle SR $+9.2$\,pp,
        collision $-2.3$\,pp, offroad $-2.2$\,pp, pedestrians flat).
\end{enumerate}

Three full-length ($3\times10^{6}$ steps) replicates of the identical
pre-campaign trunk quantify run-to-run variance: cumulative vehicle SR
ranged 52.2--62.4\% across them under the same seed --- CARLA is not
deterministic under load even when seeded. This $\sim$10\,pp band is the
yardstick against which every effect in \cref{ch:results} is read, and
motivates both the paired design and the caution attached to
single-replicate deltas (\cref{sec:threats}).

\begin{table}[t]
  \centering\small
  \caption{Condensed candidate registry (development phase, paired A/B at
  $3\times10^{5}$ steps unless noted). Deltas are versus the paired
  baseline of their era and are \emph{not} comparable across environment
  versions. Full detail in \cref{app:registry}.}
  \label{tab:candidates}
  \begin{tabular}{@{}p{4.2cm}p{2.6cm}p{2.2cm}p{3.6cm}@{}}
    \toprule
    Candidate & Type & Verdict & Headline \\
    \midrule
    Diagnostic instrumentation & measurement & adopted & no policy change \\
    Route-aware vehicle obs & observation & adopted & --- \\
    Reward shaping block (anti-stall, min-speed) & reward & promoted & core of \cref{sec:rewards} \\
    Early stuck termination & env dynamics & rejected & SR $-2.9$\,pp, s$+$t $+7.1$\,pp \\
    Route-length fix & bug fix & adopted & $\approx +50$\,pp SR (distribution change) \\
    Route-seed fix & bug fix & adopted & restores A/B pairing \\
    Value-loss repair & optimizer & retained (gate fail) & expl.\ var.\ $0 \to 0.87$; collision $+3.3$\,pp \\
    Discount $0.99 \to 0.997$ & optimizer & retained (falsified) & timeout $\to$ collision $\approx$ 1:1 \\
    Entropy schedule (fraction-based) & optimizer & promoted & stable mechanism, no instability \\
    Collision penalty $-50 \to -500$ & reward & retained (falsified) & collision flat ($-0.2$\,pp) \\
    Progress-guard removal & reward & promoted & SR $+4.7$\,pp, s$+$t $-3.5$\,pp \\
    47D observability (O1$+$O2) & observation & adopted & healthy baseline, expl.\ var.\ 0.983 \\
    Pedestrian route/pace bundle & reward $+$ routes & partial & pace out of band; vehicle ripple \\
    Pace-band retune & reward & rejected & overshoot worse; Q4 decline \\
    Hazard-gate tolerance $0.75 \to 0.85$ & reward & promoted & 5/5: SR $+2.3$\,pp, s$+$t $-3.0$\,pp, coll./off.\ improved \\
    \texttt{route\_short} as success & measurement & reverted & violated success criterion \\
    Crosswalk cap $8 \to 3$ & routes & promoted & SR $+9.2$\,pp, coll.\ $-2.3$\,pp, off.\ $-2.2$\,pp \\
    \bottomrule
  \end{tabular}
\end{table}

% ------------------------------------------------------------
\section{Final Campaign Design}
\label{sec:campaign}

The final experiment is a $2 \times 2$ factorial:
\{MLP, GNN critic encoder\} $\times$ \{curriculum, mixed-batch\} --- four
runs of $3\times10^{6}$ steps each, all on the same frozen trunk, sharing
a single experiment seed (999) paired across all four cells, with
difficulty defined by route length (\cref{sec:routes}). Run
configurations are byte-identical across cells except for the two
experimental flags (regime and critic encoder), a property verified by
direct configuration diff. The design was originally sized at three
paired seeds per cell; the seed-replication axis was dropped for
compute-budget reasons (each run occupies the single available machine
for roughly two days, \cref{app:repro}). This reduction is declared here
rather than hidden: with one seed per cell, every contrast in
\cref{ch:results} is a \emph{paired case study}, not an estimate of a
distribution.

\paragraph{Contrasts.}
The primary contrast is curriculum versus mixed-batch \emph{within}
architecture. Crucially, it is observed \emph{twice} --- once under each
critic architecture, trained independently --- so architecture
replication partially substitutes for seed replication on the primary
axis: a regime effect that appears with the same sign and comparable
structure in both pairs is much harder to attribute to run-to-run noise
than a single contrast would be. The secondary contrast is MLP versus GNN
critic \emph{within} regime; the interaction reading (hypothesis H3,
\cref{sec:rq}) is directional only at this sample size. No formal
statistical inference is claimed anywhere: effects are read against the
same-seed replicate band of \cref{sec:iterations} ($\sim$10\,pp of
vehicle SR), and corroborated by consistency across difficulty levels,
agent classes, evaluation scenarios and measurement layers rather than by
tests.

\paragraph{Measurement.}
Each run is measured at the three levels of \cref{sec:eval-pipeline}:
cumulative training metrics, static recomputation from logs, and the
final 400-episode checkpoint evaluation (three Town03 scenarios plus the
Town05 test scenario for map generalization), all reported per class.

\paragraph{Declared scope.}
The campaign does \emph{not} test: difficulty families based on traffic
density or mixed axes (defined in configuration but unused); the
attention critic encoder and PopArt normalization (available, disabled);
actor architectures (always MLP); agent-count ablations (always
$3+3$); or map pairs beyond Town03 $\to$ Town05.

The realized run matrix, with integrity checks and completion status, is
reported in \cref{tab:runmatrix}.
